% Copyright (c) 2023 Daniele Zambelli - daniele.zambelli@gmail.com

% --------------------------------
% Variabili che devono essere istanziate e che vengono usate da 
% colophon.
% \newcommand{\serie}{Matematica Libera}
% \newcommand{\titolodritto}{\volume \numero: Calcolo infinitesimale}
% \newcommand{\editore}{Pippo}
% \newcommand{\anno}{2023}
% \newcommand{\coord}{Daniele~Zambelli}
% \newcommand{\autori}{Bruno~Stecca, Daniele~Zambelli}
% \newcommand{\colab}{}
% \newcommand{\texcol}{Claudio~Beccari, Daniele~Zambelli}
% \newcommand{\docvers}{\texttt{0.9.9}}
% \newcommand{\mcisbn}{}
% \pagenumbering{roman}
\thispagestyle{empty}
\pdfbookmark{colophon}{colophon}
\vspace*{\fill}
{\setlength{\parindent}{0em}\small{
% \begin{center}
{\large{\serie\ – \titolodritto}}

Copyright {\textcopyleft} {\anno} \editore
% \end{center}

\affiancati{.70}{.28}{
Questo libro, eccetto dove diversamente specificato, è rilasciato nei 
termini della licenza Creative Commons Attribuzione – Condividi allo stesso 
modo 3.0 Italia (CC BY-SA 3.0) il cui testo integrale è disponibile agli 
indirizzi:
}{
\begin{center} \includegraphics{\magdir img/by-sa.png} \end{center}
}
\url{https://creativecommons.org/licenses/by-sa/3.0/it/}\\
\url{https://creativecommons.org/licenses/by-sa/3.0/it/legalcode}

% Sei libero di:
% riprodurre, distribuire, comunicare al pubblico, esporre in pubblico, 
% rappresentare, eseguire e recitare quest'opera, di modificare quest'opera, 
% alle seguenti condizioni:
% 
% \emph{Attribuzione} --- Devi attribuire la paternità dell'opera nei modi 
% indicati dall'autore o da chi ti ha dato l'opera in licenza e in modo tale 
% da non suggerire che essi avallino te o il modo in cui tu usi l'opera.
% 
% \emph{Condividi allo stesso modo} --- Se alteri o trasformi quest'opera, 
% o se la usi per crearne un'altra, puoi distribuire l'opera risultante solo 
% con una licenza identica o equivalente a questa.

Per maggiori informazioni su questo particolare regime di diritto d'autore 
si legga il materiale informativo pubblicato in:\\
\url{http://www.copyleft-italia.it}.

\vspace{-1em}
\paragraph{Credits:} Il testo deriva da ``Matematica Dolce'':\\ 
\url{https://www.matematicadolce.eu/} \\
che a sua volta è derivato da
``Matematica C3'':\\
\url{https://it.m.wikibooks.org/wiki/Algebra_1} \\
\url{https://it.m.wikibooks.org/wiki/Algebra_2}\\
\url{https://www.larapedia.com/geometria_razionale/
geometria_razionale_nozioni_fondamentali.html}

\vspace{-1em}
\paragraph{Coordinatori del Progetto:} \coord.

\vspace{-1em}
\paragraph{Autori:} \autori\ oltre agli autori di Matematica Dolce.

% \paragraph{Collaboratori:} \colab.

\vspace{-1em}
\paragraph{Progettazione e Implementazione in \LaTeX:} 
          {\texcol}.

\vspace{-1em}
\paragraph{Collaborazione, commenti e suggerimenti:} 
Se vuoi contribuire anche tu al progetto \serie\ o se  vuoi inviare 
commenti o suggerimenti, scrivi a:
\href{mailto:daniele.zambelli@gmail.com}{daniele.zambelli@gmail.com}.

\vspace{2ex}
Versione del documento: {\docvers} del {\oggi}.

% ISBN \mcisbn

\vspace{2ex}
{\scshape{Dati tecnici per l'adozione del libro a scuola}}

Titolo: \serie, \tipo \numero: \titolo\ - \anno.

Codice: \mcisbn 

% Editore: % \href{http://www.matematicamente.it}{\editore}. 

% Anno di edizione: \anno.

% Formato: ebook (\scshape{pdf}).
\cleardoublepage
% \pagebreak
}}
